\subsection{Orthonormal families}

DEFINITION 3. --- In a prehilbertian space, a family \((e_{i})_{i\in\mathbf{I}}\) of vectors is said to be orthogonal if \(e_{i}\) and \(e_{k}\) are orthogonals for all \(i\neq k\) , and is said to be orthonormal, if in addition \(\|e_{i}\|=1\) for all \(i\in I\) .

A subset S of E such that the family defined by the identity mapping from S onto itself is orthonormal is said to be an orthonormal set. If \((e_{i})_{i\in I}\) is an orthonormal family, the mapping \(i\mapsto e_{i}\) is injective; we can then talk indifferently of an orthonormal family or an orthonormal set.

If \((e_i)_{i\in I}\) is an orthonormal family, the complete one dimensional vector subspaces \(\mathbf{D}_i = \mathbf{K}e_i\) are two by two orthogonal. For every \(x\in E\) , the orthogonal projection of \(x\) on \(\mathbf{D}_i\) is \(\lambda_i e_i\) with \(\langle e_i|x - \lambda_i e_i\rangle = 0\) , which gives \(\langle e_i|x\rangle = \lambda_i\langle e_i|e_i\rangle = \lambda_i\) . The results of No.~2 applied to the subspaces \(\mathbf{D}_i\) imply the following propositions:

PROPOSITION 3. --- In a Hausdorff prehilbertian space E, every orthonormal family is topologically independent.

We note that this property immediately follows from the characterization of topologically independent families (IV, p.~1 and II, p.~43, cor. 2), on account of the identification of the dual of E with the completion of E or with the space conjugate to E according as K is equal to R or C (V, p.~17, Remark).

PROPOSITION 4. --- Let E be a Hausdorff prehilbertian space, \((e_{i})_{i\in I}\) an orthonormal family in E, V the closed vector subspace of E generated by the \(e_{i}\) .

\begin{enumerate}
\def\labelenumi{\arabic{enumi})}
\tightlist
\item
  For every \(x \in E\) , we have
\end{enumerate}

\[
\sum_ {i \in \mathbf {I}} \left| \langle e _ {i} | x \rangle \right| ^ {2} \leqslant \| x \| ^ {2}\tag{1}
\]

(Bessel's inequality); here the set of all \(i \in \mathbf{I}\) such that \(\langle e_i|x\rangle \neq 0\) is countable. Moreover, the following conditions are equivalent: a) \(x \in \mathbf{V}\) ; b) \(\|x\|^2 = \sum_{i \in \mathbf{I}} |\langle e_i|x\rangle|^2\) ; c) the family \(\langle e_i|x\rangle\) . \(e_i\) is summable in E, and \(x = \sum_{i \in \mathbf{I}} \langle e_i|x\rangle\) . \(e_i\) .

\begin{enumerate}
\def\labelenumi{\arabic{enumi})}
\setcounter{enumi}{1}
\item
  If \(\mathbf{V}\) is complete, then the family of all \(\langle e_i|x\rangle\) . \(e_i\) is summable in \(\mathbf{E}\) for all \(x\in\mathbf{E}\) , and \(\sum_{i\in\mathbf{I}}\langle e_i|x\rangle\) . \(e_i=p_{\mathbf{V}}(x)\) , \(\sum_{i\in\mathbf{I}}|\langle e_i|x\rangle|^2=\|p_{\mathbf{V}}(x)\|^2\) .
\item
  Suppose V is complete. For every family \((\lambda_{i})_{i\in\mathbf{I}}\) of scalars such that \(\sum_{i\in I}|\lambda_{i}|^{2}<+\infty\) , there exists a unique point \(x\in V\) such that \(\langle e_{i}|x\rangle=\lambda_{i}\) for all \(i\in I\) . If \((\mu_{i})_{i\in I}\) is a second family of scalars such that \(\sum_{i\in I}\left|\mu_{i}\right|^{2}<+\infty\) , and if \(y\in V\) is such that \(\langle e_{i}|y\rangle=\mu_{i}\) for all \(i\in I\) , then \(\langle x|y\rangle=\sum_{i\in I}\overline{\lambda}_{i}\mu_{i}\) .
\end{enumerate}

PROPOSITION 5. --- Let \((e_i)_{i \in I}\) be an orthonormal family in a Hausdorff prehilbertian space E. The following properties are equivalent :

\begin{enumerate}
\def\labelenumi{\alph{enumi})}
\item
  the family \((e_i)\) is total;
\item
  for every \(x \in E\) , the family \(\langle e_{i}|x\rangle\) . \(e_{i}\) is summable in E, and we have \(x = \sum_{i \in I} \langle e_{i}|x\rangle\) . \(e_{i}\) ;
\item
  for every \(x \in E\) ,
\end{enumerate}

\[
\left\| x \right\| ^ {2} = \sum_ {i \in \mathbf {I}} \left| \langle e _ {i} | x \rangle \right| ^ {2}\tag{2}
\]

(Parseval's relation).

When E is hilbertian these conditions are also equivalent to :

\begin{enumerate}
\def\labelenumi{\alph{enumi})}
\setcounter{enumi}{3}
\tightlist
\item
  the relations \(\langle e_i|x\rangle = 0\) for all \(i\in \mathrm{I}\) imply that \(x = 0\) .
\end{enumerate}

The equivalence of conditions a), b), c) follows immediately from prop. 4. When E is hilbertian, the equivalence of conditions a) and d) follows from cor. 1 of V, p.~16.

DEFINITION 4. --- An orthonormal and total family in a Hausdorff prehilbertian space E is called an orthonormal basis of E.

An orthonormal basis of a Hausdorff prehilbertian space E is also an orthonormal basis of the completion of E.

Let \((e_{i})_{i\in I}\) be an orthonormal basis of E; for every \(x\in E\) , the numbers \(\langle e_{i}|x\rangle\) are called, by abuse of language, the coordinates of x with respect to the basis \((e_{i})\) . For every x and y in E, we have

\[
\langle x | y \rangle = \sum_ {i \in \mathbf {I}} \overline {{\langle e _ {i} | x \rangle}} \langle e _ {i} | y \rangle .\tag{3}
\]

An orthonormal basis of E is not, in general, a basis of E over the field K in the sense defined in A, II, p.~25; to avoid any confusion we shall always say that a basis of a prehilbertian space E in the sense of loc. cit. is an algebraic basis of E over K.

Let E and F be two Hausdorff prehilbertian spaces and u a continuous linear mapping from E into F. Let \((e_{i})_{i\in I}\) (resp. \((f_{j})_{j\in J}\) ) be an orthonormal basis of E (resp. F). Put

\[
u _ {j i} = \langle f _ {j} | u (e _ {i}) \rangle
\]

for \(i \in \mathbf{I}\) , \(j \in \mathbf{J}\) . The family \((u_{ji})_{(i,j) \in \mathbf{I} \times \mathbf{J}}\) is called the matrix of \(u\) with respect to the ortho-normal bases \((e_i)\) and \((f_j)\) . Let \(x \in \mathbf{E}\) and \(y = u(x)\) ; if we write \(\xi_i = \langle e_i | x \rangle\) and \(\eta_j = \langle f_j | y \rangle\) for the coordinates of \(x\) and \(y\) respectively, we get \(\eta_j = \sum_{i \in \mathbf{I}} u_{ji} \xi_i\) for all \(j \in \mathbf{J}\) . When \((e_i)\) is an algebraic basis of \(E\) and \((f_j)\) an algebraic basis of \(F\) , our definition is consistent with that of A, II, § 10, No.~4.

Example. --- Let E be the space of all complex valued continuous functions on R, such that \(f(x + n) = f(x)\) for \(x \in R\) and \(n \in Z\) . We assign to E the scalar product defined by

\[
\langle f | g \rangle = \int_ {0} ^ {1} \overline {{{{f (t)}}}} g (t) d t.
\]

Then E is a Hausdorff prehilbertian space, but is not complete. For every integer \(n \in Z\) , let \(e_{n}(x) = \mathbf{e}(nx)\) . It is immediate that the family \((e_{n})_{n \in \mathbf{Z}}\) is orthonormal in E. Moreover, the topology of uniform convergence on E is finer than the topology deduced from the norm \(\|f\|_{2} = \langle f | f \rangle^{1/2}\) . The family \((e_{n})_{n \in \mathbf{Z}}\) is total in E for the uniform convergence (GT, X, § 4, No.~4), and a fortiori in the prehilbertian space E. Hence \((e_{n})_{n \in \mathbf{Z}}\) is an orthonormal basis of E.

